% Generated draft: requires mathematical and source review.
\documentclass[11pt]{article}
\usepackage[margin=25mm]{geometry}
\usepackage{fontspec}
\setmainfont{Latin Modern Roman}
\usepackage{amsmath,amssymb,mathtools,unicode-math}
\setmathfont{Latin Modern Math}
\usepackage{enumitem,needspace,xurl,hyperref}
\setlength{\parindent}{0pt}
\setlength{\parskip}{5pt}
\setlength{\emergencystretch}{3em}
\newcommand{\N}{\mathbb{N}}
\newcommand{\Z}{\mathbb{Z}}
\newcommand{\Q}{\mathbb{Q}}
\newcommand{\R}{\mathbb{R}}
\newcommand{\C}{\mathbb{C}}
\newcommand{\F}{\mathbb{F}}
\newcommand{\A}{\mathbb{A}}
\newcommand{\PP}{\mathbb{P}}
\newcommand{\E}{\mathbb{E}}
\newcommand{\eps}{\varepsilon}
\newcommand{\dd}{\,\mathrm{d}}
\begin{document}
\section*{UnsolvedMath ID: 30004700}
\textbf{Source problem code:} catalogue entry
\subsection{Definitions and mathematical statement}
\textbf{Scope and unresolved terminology.} The supplied record concerns reading mathematical classics across historical and cultural traditions. It does not define a mathematical classic, an alien tradition, an irretrievably lost tradition, reading, distance, otherness, or an admissible interpretive method. In particular, distance is not specified as a mathematical metric, and no mathematical domain, hypotheses, quantifiers, or criterion for a satisfactory answer are supplied.

\textbf{Source question.} What does it mean to read a mathematical classic from an alien tradition or one that is irretrievably lost, and what distance or otherness and interpretive methods does that imply?

This is a methodological and historiographical question. The supplied evidence does not determine a precise mathematical proposition; introducing formal definitions or a theorem would add unsupported content. The question is therefore retained with its unresolved qualifications for human review.
\subsection{Short English statement}
The question asks how to interpret a mathematical classic belonging to an unfamiliar or irretrievably lost tradition, and what historical or cultural distance, otherness, and interpretive methods such reading entails. The source does not specify a mathematical formalization or a criterion for resolving the question.
\subsection{Sources}
\begin{itemize}
\item Local UnsolvedMath registry, \texttt{data/unsolvedmath/problems.json}, permanent ID \(30004700\), title: Reading Mathematical Classics from Lost Traditions; source code: \texttt{catalogue entry}. The supplied record belongs to the Oberwolfach Reports Open Problems collection. No individual proposer is identified; the proposed year is recorded as 2021.
\item Supplied source citation: \emph{Mathematics and its Ancient Classics Worldwide: Translations, Appropriations, Reconstructions, Roles} (2022), Oberwolfach Reports, DOI 10.4171/owr/2021/26, pp. 1347--1406. \url{https://doi.org/10.4171/owr/2021/26}. The supplied background describes Oberwolfach Reports as documenting mathematical workshops and their open-problem discussions.
\item Supplied literature source, labelled Oberwolfach report: \url{https://doi.org/10.4171/owr/2021/26}. Literature assessment dated exactly 2026-08-22: the record labels the question open and describes it as a continuing methodological and historiographical problem with no definitive resolution. This is supplied dated curation and literature triage, not an independently verified assessment; time-sensitive status and cited claims require independent verification.
\end{itemize}
\end{document}
